\documentclass[10pt,a4paper]{amsart}
\usepackage[margin=19mm]{geometry}
\usepackage{amsmath,amssymb,mathtools}
\usepackage[scr=rsfs,cal=euler]{mathalfa}
\usepackage{xfrac}
\usepackage[unicode,hidelinks,bookmarks=false]{hyperref}
\newcommand{\ie}{i.e.\ }
\newcommand{\cf}{cf.\ }
\newcommand{\resp}{respectively\ }
\newcommand{\Subsection}{\subsection}
\providecommand{\nobreakdash}{\nobreak\hbox{-}}
\setlength{\emergencystretch}{2em}
\multlinegap0pt
\usepackage{bm}
\usepackage{bbm}
\usepackage{dsfont}
\usepackage{xspace}
\usepackage{cancel}
\usepackage{mathtools}
\usepackage{xcolor}
\usepackage{amsthm}
\makeatletter
\@ifundefined{proposition}{\newtheorem{proposition}{Proposition}}{}
\@ifundefined{theorem}{\newtheorem{theorem}{Theorem}}{}
\makeatother
\DeclareMathOperator{\Reg}{Reg}
\newcommand{\ch}[1]{{\color{black} #1}}
\newcommand{\mb}{\mathbb}
\title[Recovery problem of parametrizations from Legendre data]{Recovery problem of parametrizations from Legendre data}
\author{C. Mu\~noz-Cabello, T. Nishimura, R. Oset Sinha}
\date{Source: arXiv:2602.19924v1 (CC BY 4.0)}
\begin{document}
\maketitle

\noindent\emph{Source and licence.} Recovery problem of parametrizations from Legendre data. Version arXiv:2602.19924v1 (CC BY 4.0), distributed under the Creative Commons Attribution 4.0 International licence, as recorded in the arXiv licence metadata for this exact version and shown on the abstract page https://arxiv.org/abs/2602.19924v1. Full rights notice: licenses/arxiv-2602.19924-CC-BY-4.0.txt.\par\smallskip

\noindent\textit{Editorial scope.} Counted unit: the Theorem whose source label is \texttt{theorem1} in the article (source0.tex lines 748--762, source label \texttt{theorem1}), delivered together with the article's own complete original proof of it (source0.tex lines 928--1026, one \texttt{proof} environment of 99 lines). No mathematics is written, repaired or re-derived: statement and proof are the article's own text line for line. The proof is detached from the statement in the source: the article states the result at lines 748--762 and prints its proof at lines 928--1026, and the proof environment's own header names the statement, so the association is the article's own. Required context retained uncounted: nu mu orthonormal basis (lines 920--926). Every definition, display and label of those lines is retained unchanged. Omitted from this package and not redistributed: 1--747, 763--919, 927--927, 1027--2464. Nothing inside a retained excerpt is omitted or altered: every retained body is delivered exactly as the article's lines are written, so no omission receipt of any kind accompanies this package. Citations: the retained excerpts carry no citation command at all, so no bibliography entry of the article is transcribed and no reference is redistributed. Reference adaptation: none was needed and none was made; every reference inside the delivered unit resolves inside it, so no unresolved reference or citation is produced. Printed slips kept literal: no printed word, symbol, hypothesis or number of the counted statement or of its proof was altered. Layout and class: the article's own class is \texttt{amsart} with packages including amscd, amsfonts, amsmath, amssymb, latexsym, amsthm, color, xy, epsfig, graphicx, todonotes; the wrapper is the lane's pinned \texttt{amsart} editorial wrapper at 10pt on A4, which restates the theorem-like environments the retained text needs and adds the article's own macro definitions that the retained text needs (\texttt{\textbackslash Reg}, \texttt{\textbackslash ch}, \texttt{\textbackslash mb}), with extra wrapper packages (bm, bbm, dsfont, xspace, cancel, mathtools, xcolor). The delivered numbering is the unit's own. Assets: no retained span contains a figure environment, a graphics-inclusion command, a PDF-page inclusion, a table or a TikZ picture; no image file is copied into this package, the package declares no asset dependency and no image file is redistributed with it. Licence: version arXiv:2602.19924v1 of this article is distributed under the Creative Commons Attribution 4.0 International licence, as recorded in the arXiv licence metadata for this exact version and shown on the abstract page https://arxiv.org/abs/2602.19924v1. A full rights notice is delivered at licenses/arxiv-2602.19924-CC-BY-4.0.txt. Characters: every retained excerpt is pure ASCII, so the delivered engine, XeLaTeX with its default Latin Modern text font, reports no missing character.
\begin{proposition}\label{nu mu orthonormal basis}
	The set $\{\nu^n({\color{black}\boldsymbol{\theta}}),
	\hat\mu_1({\color{black}\boldsymbol{\theta}}),
	\dots,\hat\mu_n({\color{black}\boldsymbol{\theta}})\}$ 
	is an orthonormal basis for $T_{\nu({\color{black}\boldsymbol{\theta}})}
	\mathbb{R}^{n+1}$, where $\hat\mu_i=\tilde\mu_i/|\tilde\mu_i|$ for $i=1,\dots,n$.
\end{proposition}

\begin{theorem}\label{theorem1}
Let $f: U_n\to \mathbb{R}^{n+1}$ be a {regular} frontal.   
Then, $f$ is recovered from its Legendre data 
$\left\{\nu({\color{black}\bf x}), 
a({\color{black}\bf x})\right\}_{{\color{black}\bf x}\in Reg({\color{black}\nu})}$.   
{\color{black} 
Namely, for any ${\color{black}\bf x}\in U_n$, 
$f({\color{black}\bf x})$ can be concretely described in terms of 
$\nu\left({\color{black}\bf x}_i\right), 
a\left({\color{black}\bf x}_i\right)$ and $\lim_{i\to \infty}$ where 
$\left\{{\color{black}\bf x}_i\ldots\right\}_{i=1, 2, \ldots}$ 
is any sequence of regular points of $\nu: U_n\to S^n$ 
such that $\lim_{i\to \infty}{\color{black}\bf x}_i={\color{black}\bf x}$.     
}
\end{theorem}

\begin{proof}[Proof of Theorem \ref{theorem1}]
	Given ${\color{black}\bf x} \in U_n$, 
	we identify $\mb{R}^{n+1}$ with $T_{f({\color{black}\bf x})}
	\mb{R}^{n+1}$.
	Let $\nu\colon \Reg(f) \to S^n$ be the unit vector field along $f$ and $a=f\cdot \nu$.	
	Since $|\nu({\color{black}\bf x})|=1$ 
	for all ${\color{black}\bf x} \in U_n$, we can choose $\theta_1,\dots,\theta_n\colon \Reg(f) \to W$ such that
		\[\nu=\nu\circ (\theta_1,\dots,\theta_n).\]
	Moreover, $\nu$ verifies the identity $\nu({\color{black}\bf x})
	\cdot df_{\color{black}\bf x}({\color{black}\bf v})=0$ 
	for all ${\color{black}\bf v} \in T_{\color{black}\bf x}\mb{R}^n$ and 
	${\color{black}\bf x} \in \Reg(f)$, hence
		\[\frac{\partial a}{\partial x_i}({\color{black}\bf x})
		=\frac{\partial f}{\partial x_i}({\color{black}\bf x})
		\cdot \nu({\color{black}\bf x})+f({\color{black}\bf x})
		\cdot \frac{\partial \nu}{\partial x_i}({\color{black}\bf x})
		=f({\color{black}\bf x})\cdot 
		\frac{\partial \nu}{\partial x_i}({\color{black}\bf x}).\]
	Using the chain rule, we obtain the following system of differential equations:
	\[
		\frac{\partial a}{\partial x_i}=f\cdot \frac{\partial \nu^n}{\partial x_i}
		=f\cdot \frac{\partial \nu^n}{\partial \theta_1}\frac{\partial \theta_1}{\partial x_i}+\dots+f\cdot 	\frac{\partial \nu^n}{\partial \theta_n}\frac{\partial \theta_n}{\partial x_i}=b_1\frac{\partial \theta_1}{\partial x_i}+\dots+b_n\frac{\partial \theta_n}{\partial x_i},
	\]
	which can be written in matrix form as
	\begin{align}\label{system b1 bn}
		\begin{pmatrix}
			\dfrac{\partial a}{\partial x_1} \\ \vdots \\ \dfrac{\partial a}{\partial x_n}	
		\end{pmatrix}
		=
		\begin{pmatrix}
			\dfrac{\partial \theta_1}{\partial x_1} & \hdots & \dfrac{\partial \theta_n}{\partial x_1}	\\
			\vdots				& \ddots & \vdots				\\
			\dfrac{\partial \theta_1}{\partial x_n} & \hdots & \dfrac{\partial \theta_n}{\partial x_n}
		\end{pmatrix}
		\begin{pmatrix}
			b_1 \\ \vdots \\ b_n
		\end{pmatrix}
	\end{align}

	{\color{black}Notice that} the mapping 
	$d\nu_{\color{black}\bf x}\colon T_{\color{black}\bf x}\Reg(f) 
	\to T_{\nu({\color{black}\bf x})}S^n$ 
	is a monomorphism {\color{black}for any ${\bf x} \in \Reg(\nu)$}.    
	By the chain rule, 
		\[d\nu_{\color{black}\bf x}=
		d\nu^n_{\theta({\color{black}\bf x})}\circ d\theta_{\color{black}\bf x},\]
	hence $d\theta_{\color{black}\bf x}\colon 
	T_{\color{black}\bf x}\Reg(\nu) \to 
	T_{\theta({\color{black}\bf x})}W$ is also a monomorphism.
	Since $\dim T_{\color{black}\bf x}\Reg(\nu) 
	= \dim T_{\theta({\color{black}\bf x})}W$, 
	the Grassman formula implies that $\operatorname{rk} 
	d\theta_{\color{black}\bf x}=\dim T_{\color{black}\bf x}U$, 
	so $d\theta_{\color{black}\bf x}$ is an isomorphism and the coefficient matrix of Equation \eqref{system b1 bn} is invertible at ${\color{black}\bf x}$.
	It follows that the system of equations has a unique solution $b_1,\dots,b_n\colon \Reg(\nu) \to \mb{R}$.

	If we set $\mu_i=\hat\mu_i\circ \theta$ for $i=1,\dots,n$, it follows from Proposition \ref{nu mu orthonormal basis} that we can write 
		\[
		f({\color{black}\bf x})
		=f({\color{black}\bf x})\cdot 
		\nu({\color{black}\bf x})\nu({\color{black}\bf x})
		+f({\color{black}\bf x})\cdot 
		\mu_1({\color{black}\bf x})\mu_1({\color{black}\bf x})
		+\dots+f({\color{black}\bf x})\cdot 
		\mu_n({\color{black}\bf x}) \mu_n({\color{black}\bf x}).
		\] 
	On the one hand, 
	we know that $f({\color{black}\bf x})\cdot 
	\nu({\color{black}\bf x})=
	a({\color{black}\bf x})$ by definition of height function.
	On the other hand, for $j=1,\dots,n$,
		\[f({\color{black}\bf x})\cdot\mu_j({\color{black}\bf x})
		=f({\color{black}\bf x})\cdot \hat 
		\mu_j(\theta({\color{black}\bf x}))
		=f({\color{black}\bf x})\cdot 
		\frac{\tilde \mu_j(\theta({\color{black}\bf x}))}
		{|\tilde \mu_j(\theta({\color{black}\bf x}))|}
		=\frac{f({\color{black}\bf x})}
		{|\tilde \mu_j(\theta({\color{black}\bf x}))|}
		\cdot\frac{\partial \nu^n}{\partial \theta_j}
		(\theta({\color{black}\bf x}))
		=\frac{b_j({\color{black}\bf x})}
		{|\tilde\mu_j(\theta({\color{black}\bf x}))|}.\]
	Therefore,
	\begin{equation}\label{equation to recover f}
		f({\color{black}\bf x})
		=a({\color{black}\bf x})\nu({\color{black}\bf x})
		+\frac{b_1({\color{black}\bf x})}{|\tilde\mu_1(\theta({\color{black}\bf x}))|}
		\mu_1({\color{black}\bf x})+\dots
		+\frac{b_n({\color{black}\bf x})}{|\tilde\mu_n(\theta({\color{black}\bf x}))|}
		\mu_n({\color{black}\bf x})
	\end{equation}
	for any ${\color{black}\bf x} \in \Reg(\nu)$.
	Note that every function on the right-hand side of this identity can be computed entirely in terms of $a$ and $\nu$.
	
	The set $\Reg(\nu)$ is \ch{open and} dense in $\Reg(f)$ by definition of pseudo regular frontal, and since $\Reg(f)$ is \ch{open and} dense in $U_n$, $\Reg(\nu)$ is \ch{open and} dense in $U_n$.
	\text{Since they are continuous}, we can uniquely extend the mapping $\theta\colon \Reg(f) \to W$ and the functions $b_1,\dots,b_n\colon \Reg(\nu) \to \mb{R}$ to $U_n$, which is the domain of $f$.
	We conclude that the identity \eqref{equation to recover f} holds in $U_n$, so it can be written entirely in terms of the Legendre data, as stated.
\end{proof}

\end{document}
